% Generated draft: requires mathematical and source review.
\documentclass[11pt]{article}
\usepackage[margin=25mm]{geometry}
\usepackage{fontspec}
\setmainfont{Latin Modern Roman}
\usepackage{amsmath,amssymb,mathtools,unicode-math}
\setmathfont{Latin Modern Math}
\usepackage{enumitem,needspace,xurl,hyperref}
\setlength{\parindent}{0pt}
\setlength{\parskip}{5pt}
\setlength{\emergencystretch}{3em}
\newcommand{\N}{\mathbb{N}}
\newcommand{\Z}{\mathbb{Z}}
\newcommand{\Q}{\mathbb{Q}}
\newcommand{\R}{\mathbb{R}}
\newcommand{\C}{\mathbb{C}}
\newcommand{\F}{\mathbb{F}}
\newcommand{\A}{\mathbb{A}}
\newcommand{\PP}{\mathbb{P}}
\newcommand{\E}{\mathbb{E}}
\newcommand{\eps}{\varepsilon}
\newcommand{\dd}{\,\mathrm{d}}
\begin{document}
\section*{UnsolvedMath ID: 8800023}
\textbf{Source problem code:} catalogue entry
\subsection{Definitions and mathematical statement}
Let $N\geq4$ be composite and let
\[
R(N)=\{\overline x\in\mathbb Z/N\mathbb Z:x^2\equiv1\pmod N\}.
\]
Develop an efficient algorithm for, or a characterization of, the subset of relevant roots intended in the slides. The extract does not define relevant, give a small-root bound or interval, or state any auxiliary information about $N$.

The classes $\pm1$ always satisfy the congruence, so finding an arbitrary root would not express a nontrivial task. A nontrivial root means one not congruent to either $1$ or $-1$. In the classical special case $N=pq$ with distinct odd primes, such a root yields factors through $\gcd(x-1,N)$ and $\gcd(x+1,N)$, as reported in the supplied summary. This fact does not supply an efficient way to find the root from $N$ alone or identify the missing small-root variant.

The general composite input is not silently restricted to two primes. A complete algorithmic target must specify root representatives and size bounds, what factorization information is supplied, and whether output is one selected root, an explicit enumeration, or a compressed characterization. If a runtime bound is in input bit length, write it in terms of $\log N$ and account for output size when listing all roots. The source gives no exact efficiency or success-probability requirement.
\subsection{Short English statement}
Find or characterize the roots of $x^2\equiv1\pmod N$ selected by the missing slide conditions. Trivial roots and the classical factoring implication do not resolve the unspecified small-root task.
\subsection{Sources}
\begin{itemize}
\item Alexander May, IPAM \emph{Workshop on Number Theory and Cryptography --- Open Problems} (2006), attachment 6222, PDF page 17, printed slide 12, schedule source \url{https://www.ipam.ucla.edu/programs/workshops/workshop-i-number-theory-and-cryptography-open-problems/?tab=schedule}; accessed 2026-07-29. The supplied provenance records a PDF/legacy-PPT visual audit.
\item Local registry \texttt{data/unsolvedmath/problems.json}, permanent UnsolvedMath ID $8800023$, source code catalogue entry, \emph{Roots of x-squared minus one modulo a composite}. Dataset attribution: Ulam AI and the UnsolvedMath Contributors.
\item The supplied summary mentions the classical nontrivial-square-root factoring step and Miller--Rabin/splitter context without a separate publication citation.
\end{itemize}
\end{document}
